\documentclass[10pt,a4paper]{amsart}
\usepackage[margin=19mm]{geometry}
\usepackage{amsmath,amssymb,mathtools}
\usepackage[scr=rsfs,cal=euler]{mathalfa}
\usepackage{xfrac}
\usepackage[unicode,hidelinks,bookmarks=false]{hyperref}
\newcommand{\ie}{i.e.\ }
\newcommand{\cf}{cf.\ }
\newcommand{\resp}{respectively\ }
\newcommand{\Subsection}{\subsection}
\providecommand{\nobreakdash}{\nobreak\hbox{-}}
\setlength{\emergencystretch}{2em}
\multlinegap0pt
\usepackage{bm}
\usepackage{bbm}
\usepackage{dsfont}
\usepackage{xspace}
\usepackage{cancel}
\usepackage{mathtools}
\usepackage{amsthm}
\makeatletter
\@ifundefined{lemma}{\newtheorem{lemma}{Lemma}}{}
\makeatother
\newcommand{\Cov}{\operatorname{Cov}}
\newcommand{\E}{{\mathbb E}}
\newcommand{\N}{{\mathbb N}}
\newcommand{\Var}{\operatorname{Var}}
\newcommand{\Z}{{\mathbb Z}}
\renewcommand{\P}{{\mathbb P}}
\renewcommand{\S}{\mathbb{S}}
\title[Furstenberg systems of pretentious and MRT multiplicative fu]{Furstenberg systems of pretentious and MRT multiplicative functions}
\author{Nikos Frantzikinakis, Mariusz Lema\'nczyk, Thierry de la Rue}
\date{Source: arXiv:2304.03121v4 (CC BY 4.0)}
\begin{document}
\maketitle

\noindent\emph{Source and licence.} Furstenberg systems of pretentious and MRT multiplicative functions. Version arXiv:2304.03121v4 (CC BY 4.0), distributed under the Creative Commons Attribution 4.0 International licence, as recorded in the arXiv licence metadata for this exact version and shown on the abstract page https://arxiv.org/abs/2304.03121v4. Full rights notice: licenses/arxiv-2304.03121-CC-BY-4.0.txt.\par\smallskip

\noindent\textit{Editorial scope.} Counted unit: the Lemma whose source label is \texttt{L:thetap'} in the article (source0.tex lines 1935--1949, source label \texttt{L:thetap'}), delivered together with the article's own complete original proof of it (source0.tex lines 1950--2080, one \texttt{proof} environment of 131 lines). No mathematics is written, repaired or re-derived: statement and proof are the article's own text line for line. Required context retained uncounted: none. Every definition, display and label of those lines is retained unchanged. Omitted from this package and not redistributed: 1--1934, 2081--4520. Nothing inside a retained excerpt is omitted or altered: every retained body is delivered exactly as the article's lines are written, so no omission receipt of any kind accompanies this package. Citations: the retained excerpts carry no citation command at all, so no bibliography entry of the article is transcribed and no reference is redistributed. Reference adaptation: none was needed and none was made; every reference inside the delivered unit resolves inside it, so no unresolved reference or citation is produced. Printed slips kept literal: no printed word, symbol, hypothesis or number of the counted statement or of its proof was altered. Layout and class: the article's own class is \texttt{amsart} with packages including hyperref, amssymb, amsfonts, textcomp, graphicx, enumerate, inputenc, fontenc, eucal, xcolor, geometry; the wrapper is the lane's pinned \texttt{amsart} editorial wrapper at 10pt on A4, which restates the theorem-like environments the retained text needs and adds the article's own macro definitions that the retained text needs (\texttt{\textbackslash Cov}, \texttt{\textbackslash E}, \texttt{\textbackslash N}, \texttt{\textbackslash P}, \texttt{\textbackslash S}, \texttt{\textbackslash Var}, \texttt{\textbackslash Z}), with extra wrapper packages (bm, bbm, dsfont, xspace, cancel, mathtools). The delivered numbering is the unit's own. Assets: no retained span contains a figure environment, a graphics-inclusion command, a PDF-page inclusion, a table or a TikZ picture; no image file is copied into this package, the package declares no asset dependency and no image file is redistributed with it. Licence: version arXiv:2304.03121v4 of this article is distributed under the Creative Commons Attribution 4.0 International licence, as recorded in the arXiv licence metadata for this exact version and shown on the abstract page https://arxiv.org/abs/2304.03121v4. A full rights notice is delivered at licenses/arxiv-2304.03121-CC-BY-4.0.txt. Characters: every retained excerpt is pure ASCII, so the delivered engine, XeLaTeX with its default Latin Modern text font, reports no missing character.
	\begin{lemma}\label{L:thetap'}
	Let  $f\colon \N\to \S^1$ be a multiplicative function  such that for every $p\in\P$ we have   $f(p^s):=e(\theta_p)$ for some $\theta_p\in [-1/2,1/2)$ and  all $s\in \N$.  For $N\in \N$, let
	\begin{equation}\label{E:AN}
		A(N):=\sum_{p\in \P\cap [N]}  \frac{\theta_p}{p}.
	\end{equation}
	Then, for some  universal constant $C>0$,  we have
	%%5the following implication holds
	%%	 	$$
	%%	\sum_{p\in\P}\frac{\theta_p^2}{p}\leq c\, \varepsilon^3 \implies
	%%	\limsup_{N\to\infty}\E_{n\in [N]}|f(n)-e(a(N))|\leq  \varepsilon.
	%%	$$
	\begin{equation}\label{E:ANest}
	\E_{n\in [N]}|f(n)-e(A(N))|^2\leq  C \,	\Big(\sum_{p\in \P\cap [N]} \frac{\theta_p^2}{p}+ \frac{\log\log{N}}{\log{N}}\Big).
	\end{equation}
\end{lemma}

\begin{proof}
For convenience, we use probabilistic language. The reader can translate the argument to conventional number theoretic language by replacing $X_p$ with $\theta_p\cdot {\bf 1}_{p\Z}$
and $\E_N$ with $\E_{n\in[N]}$ throughout.


	For a fixed $N\in \N$, we  consider the finite probability space that consists of the  interval $[N]$ together with the probability measure that  assigns mass $1/N$ to each point in $[N]$. For a prime $p\in [N]$, we define the random variable
	$ X_p\colon [N]\to \{0,\theta_p\}$ by
	$$
	X_p(n):=
	\begin{cases} \theta_p \quad & \text {if } \, p\mid n, \\
		0 & \text{otherwise}.
	\end{cases}
	$$
	These random variables are relevant for our problem because if  we let
	$$
	S_N:=\sum_{p\in \P\cap [N]}\, X_p,
	$$
	then from the definition of $f$ on powers of primes and its multiplicativity, we have
	$$
	f(n)=e(S_N(n)), \quad n\in [N].
	$$
	%%Our next goal is to estimate the variance of $S_N$, namely for
	%%$$
	%%\Var_N(S_N)=\E_{n\in[N]}|S_N(n)-a(N)|^2, \quad a(N):=\E_N(S_N).
	%%$$
	%%Since  $f(n)=e(S_N(n))$,  the last quantity is linked to the one we want to bound by the %%elementary estimate
	Note that if
	$$
	a(N):=\E_N(S_N),
	$$
	then
	\begin{equation}\label{E:Var1}
		\E_{n\in [N]}|f(n)-e(a(N))|^2\leq 4\pi^2\,  \E_{n\in [N]}|S_N(n)-a(N)|^2=4 \pi^2\cdot \Var_N(S_N),
	\end{equation}
	where to get the first estimate we used that $|e(x)-e(y)|=2\pi \big|\int_x^y e(t)\, dt\big|\leq 2\pi|x-y|$.
	
	So, our problem reduces to getting an upper  bound  for   $ \Var_N(S_N)$.
	We start with some   easy computations that give:
	%% $$
	%%\E_N(X_p)=\frac{1}{N}\, \Big\lfloor\frac{N}{p}\Big\rfloor
	%%$$
	%%$$
	%%\Var_N(X_p)= \E_N(X_p^2)-(\E_N(X_p))^2\leq \frac{1}{p}
	%%$$
	%%$$
	%%\Cov(X_p,X_q)=\E_N(X_p\cdot X_q)-\E_N(X_p)\cdot \E_N(X_q)\leq %%\frac{1}{N}\Big(\frac{1}{p}+\frac{1}{q}\Big).
	%%$$
	%%For  prime $p\in [N]$ we also define the random variables
	%%$$
	%%Y_p:=\theta_p\cdot  X_p.
	%%$$
	%%Then
\begin{equation}\label{srednia}
	\E_N(X_p)= \frac{\theta_p}{N}\, \Big\lfloor\frac{N}{p}\Big\rfloor,
\end{equation}
	$$
	\Var_N(X_p)= \E_N(X_p^2)-(\E_N(X_p))^2\leq \frac{\theta_p^2}{p},
	$$
	$$
	\Cov_N(X_p,X_q)=\E_N(X_p\cdot X_q)-\E_N(X_p)\cdot \E_N(X_q) = \frac{\theta_p\, \theta_q}{N}\left(\Big\lfloor\frac{N}{pq}\Big\rfloor - \frac{1}{N}\Big\lfloor\frac{N}{p}\Big\rfloor \cdot \Big\lfloor\frac{N}{q}\Big\rfloor\right).
	%\leq \frac{\theta_p\, \theta_q}{N}\Big(\frac{1}{p}+\frac{1}{q}\Big),
	$$
	To get the last equality we assume that $p\neq q$ and uses  that $\E_N(X_p\cdot X_q)=\frac{\theta_p\, \theta_q}{N}\, \Big\lfloor\frac{N}{pq}\Big\rfloor$, which follows from
	the primality of $p,q$.
	
	Now let us assume first  that all $\theta_p$'s have the same sign. Then using the inequality
	$$
	\Big\lfloor\frac{N}{pq}\Big\rfloor - \frac{1}{N}\Big\lfloor\frac{N}{p}\Big\rfloor \cdot \Big\lfloor\frac{N}{q}\Big\rfloor \leq \frac{1}{p}+\frac{1}{q},
	$$
	we get for all $p\neq q$
	$$
	\Cov_N(X_p,X_q) \leq \frac{\theta_p\, \theta_q}{N}\Big(\frac{1}{p}+\frac{1}{q}\Big).
	$$
	
	Using these estimates and the identity
	$$
	\Var_N(S_N)=\sum_{p\in \P\cap[N]} \Var_N(X_p)+\sum_{p,q\in \P\cap [N], p\neq q} \Cov_N(X_p\cdot X_q),
	$$
	we deduce that
	$$
	\Var_N(S_N)\leq  \sum_{p\in \P\cap[N]} \frac{\theta_p^2}{p}+
	\frac{1}{N}\sum_{p,q\in \P\cap [N], p\neq q} \theta_p\, \theta_q\Big(\frac{1}{p}+\frac{1}{q}\Big).
	$$
	Since $\theta_p\in [-1/2,1/2)$, $p\in \P$, we can bound the second sum by $\sum_{p,q\in \P\cap [N], p\neq q}\frac{1}{p}$, which in turn is bounded by $C_1 \, \frac{N}{\log{N}}\,  \log\log{N}$ for some universal constant $C_1$. Hence,
	\begin{equation}\label{E:Var2}
		\Var_N(S_N)\leq  \sum_{p\in \P\cap[N]} \frac{\theta_p^2}{p}+C_1\, \frac{\log\log{N}}{\log{N}}.
	\end{equation}
	The bound \eqref{E:ANest}, with $a(N)$  in place of $A(N)$,  then follows by combining \eqref{E:Var1} and  \eqref{E:Var2}.
	To get  \eqref{E:ANest}, as stated,  note that by~\eqref{srednia}
%	\begin{equation}\label{E:aN}
$$
		a(N)=\sum_{p\in \P\cap [N]}  \frac{\theta_p}{N}\, \Big\lfloor\frac{N}{p}\Big\rfloor
$$
%	\end{equation}
	and
	$$
	|a(N)-A(N)|\leq \frac{C_2}{\log{N}}
	$$
	for some universal constant $C_2>0$.
	
	This completes the proof in the case where all the $\theta_p$'s have the same sign.
	
	In the general case, we can decompose $f$ as a product $f_+\cdot f_-$, where $f_+$ and $f_-$ are multiplicative functions defined by
	$$ f_+(p^s) := \begin{cases}
                    e(\theta_p),\quad  & \text{if }\theta_p\geq0\\
	                1,\quad   &\text{otherwise}
	              \end{cases}
    \quad\text{ and }\quad
    f_-(p^s) := \begin{cases}
                    e(\theta_p),\quad  &\text{if }\theta_p<0\\
	                1, \quad  &\text{otherwise,}
	              \end{cases}
    $$
	for every $p\in\P$ and every $s\in\N$. Now, let
	$$
	A_+(N):=\sum_{\substack{p\in \P\cap [N],\\ \theta_p\geq0}}  \frac{\theta_p}{p}
	\quad\text{ and }\quad
	A_-(N):=\sum_{\substack{p\in \P\cap [N],\\ \theta_p<0}}  \frac{\theta_p}{p}.
	$$
	Then, we have for all $n\in[N]$
	\begin{align*}
	|f(n)-e(A(N))| &= |f_+(n)\cdot f_-(n)-e(A_+(N))\cdot e(A_-(N))| \\
	&\leq |f_+(n)-e(A_+(N))| + |f_-(n)-e(A_-(N))|.
	\end{align*}
	It follows   that
	$$
	\E_N |f(n)-e(A(N))|^2 \leq 2\cdot \left(\E_N |f_+(n)-e(A_+(N))|^2 +
	\E_N |f_-(n)-e(A_-(N))|^2\right),
	$$
	and we get the conclusion by applying the preceding analysis to $f_+$ and $f_-$.
\end{proof}	 	

\end{document}
